\documentclass[a4paper,12pt]{article}
\usepackage[utf8]{inputenc}
\usepackage[catalan]{babel}
\usepackage{amsmath}
\usepackage{amsfonts}
\usepackage{amssymb}
\usepackage{geometry}
\geometry{margin=2cm}

\title{Examen de Matemàtiques Resolt - 4t ESO\\ \large Tema: Exponencials i Logaritmes}
\author{}
\date{}

\begin{document}

\maketitle

\begin{tabular}{ll}
\textbf{Nom i cognoms:} \rule{8cm}{0.1pt} & \textbf{Grup:} \rule{2cm}{0.1pt} \\
\end{tabular}

\vspace{0.5cm}
\textit{Durada aproximada: 50 minuts}
\vspace{0.5cm}

\section*{1. Definició de logaritme [2.5 Pts]}
Troba el valor de $x$ en les igualtats següents, aplicant la definició $\log_a(b) = c \iff a^c = b$:

\begin{itemize}
    \item[a)] $\log_{2}(x)=5 \implies x = 2^5 \implies \mathbf{x = 32}$
    \item[b)] $3^{x}=81 \implies x = \log_3(81) \implies x = 4 \implies \mathbf{x = 4}$
    \item[c)] $\log_{x}(125)=3 \implies x^3 = 125 \implies x = \sqrt[3]{125} \implies \mathbf{x = 5}$
    \item[d)] $10^{x}=0.0001 \implies x = \log(0.0001) \implies x = \log(10^{-4}) \implies \mathbf{x = -4}$
    \item[e)] $2^{x+1}=1024 \implies x+1 = \log_2(1024) \implies x+1 = 10 \implies \mathbf{x = 9}$
\end{itemize}

\section*{2. Expressions [2.5 Pts]}
Simplifica al màxim les expressions següents utilitzant les propietats de les potències:

\begin{itemize}
    \item[a)] $A \cdot \sqrt{\left(\frac{A^{2} \cdot \sqrt{A}}{\sqrt{A} \cdot A}\right)}$
    
    \textbf{Solució pas a pas:}
    1. Passem totes les arrels a potències:
    \[ A^1 \cdot \sqrt{\frac{A^2 \cdot A^{1/2}}{A^{1/2} \cdot A^1}} \]
    2. Sumem els exponents del numerador i del denominador de la fracció interior:
    \[ A^1 \cdot \sqrt{\frac{A^{2 + 1/2}}{A^{1/2 + 1}}} = A^1 \cdot \sqrt{\frac{A^{5/2}}{A^{3/2}}} \]
    3. Restem els exponents (numerador menys denominador):
    \[ A^1 \cdot \sqrt{A^{5/2 - 3/2}} = A^1 \cdot \sqrt{A^{2/2}} = A^1 \cdot \sqrt{A^1} \]
    4. Multiplicació final sumant els exponents:
    \[ A^1 \cdot A^{1/2} = A^{1 + 1/2} = \mathbf{A^{3/2}} \text{ o } \mathbf{\sqrt{A^3}} \]

    \item[b)] $2 \log(2x) + \frac{1}{2} \log(4x^{6}) - \log(x)$
    
    \textbf{Solució:}
    \[ \log((2x)^2) + \log((4x^6)^{1/2}) - \log(x) = \log(4x^2) + \log(2x^3) - \log(x) \]
    \[ \log \left( \frac{4x^2 \cdot 2x^3}{x} \right) = \log \left( \frac{8x^5}{x} \right) = \mathbf{\log(8x^4)} \]
\end{itemize}

\section*{3. Propietats dels logaritmes [2 Pts]}
Sabent que $\log k = 1.2$, calcula:

\begin{itemize}
    \item[a)] $\log(100 \cdot k) = \log(100) + \log(k) = 2 + 1.2 = \mathbf{3.2}$
    \item[b)] $\log\left(\frac{k^3}{10}\right) = 3 \log(k) - \log(10) = 3(1.2) - 1 = \mathbf{2.6}$
    \item[c)] $\log(\sqrt[4]{k}) = \frac{1}{4} \log(k) = \frac{1.2}{4} = \mathbf{0.3}$
    \item[d)] $k = 10^{\log k} = 10^{1.2} \implies \mathbf{k \approx 15.85}$
\end{itemize}

\section*{4. Problema d'aplicació [3 Pts]}
Població actual: $12.500$. Creixement: $3,5\%$ anual.

\begin{equation}
    100\% + 3,5\% = 103,5\% \implies\frac{103,5}{100}=1,035
\end{equation}

\begin{itemize}
    \item[a)] $\mathbf{P(t) = 12500 \cdot 1.035^t}$
    \item[b)] $P(8) = 12500 \cdot 1.035^8 \approx \mathbf{16460 \text{ habitants}}$
    \item[c)] $P(-2) = 12500 \cdot 1.035^{-2} \approx \mathbf{11669 \text{ habitants}}$
    \item[d)] $20000 = 12500 \cdot 1.035^t \implies 1.6 = 1.035^t \implies t = \frac{\log(1.6)}{\log(1.035)} \approx \mathbf{13.66 \text{ anys}}$
\end{itemize}

\end{document}